\documentclass[12pt]{article}
\usepackage[utf8]{inputenc}
\usepackage[T1]{fontenc}
\usepackage{hyperref}
\title{Matbat Language Resources for K-12 Curriculum Development}
\author{Compiled by Claude Sonnet 5.5 (claude-sonnet-5-5) for Andr\'as Kornai}
\date{}
\begin{document}
\maketitle

\section{Introduction}
Matbat (ISO 639-3 xmt) is an Austronesian language of Misool in the Raja Ampat Islands, Indonesia. Sources give about 1,000--1,500 speakers. It is one of a small number of Austronesian languages with lexical tone, like the neighbouring Ma'ya. The search found documentation by linguists but no school materials.

\section{Primary Educational Resources}
No primers, textbooks or school programmes for Matbat were found.

\section{Linguistic and Reference Materials}
\begin{itemize}
\item Bert Remijsen's lexicon of the Magey dialect, with more than 1000 entries, collected during fieldwork in 2003 and hosted by the University of Edinburgh (DOI 10.7488/ds/262). \url{https://datashare.ed.ac.uk/handle/10283/796}
\item Lexical comparanda for 14 Austronesian varieties of Raja Ampat, a Zenodo dataset. \url{https://zenodo.org/records/14716166}
\item Overview of the language and its dialects. \url{https://en.wikipedia.org/wiki/Matbat_language}
\item A short Matbat verb list on Verbix. \url{https://docs.verbix.com/Languages/Matbat}
\item An ASJP wordlist page for the neighbouring Misool Ma'ya, for comparison only. \url{https://asjp.clld.org/languages/MISOOL_MAYA}
\end{itemize}

\section{Religious and Community Resources}
A Raja Ampat marine conservation authority page on local culture and heritage appeared in the results. It was not checked for Matbat content. \url{https://kkprajaampat.com/en/culture-heritage/} No Matbat church, scripture or community organisation was identified. A page on the related Matlow people is at \url{https://en.wikipedia.org/wiki/Matlow_people}.

\section{Digital and Audio Resources}
No Matbat audio, video or software was found in the searches. The Edinburgh dataset may include audio, but that was not confirmed.

\section{Curriculum Development Context}
The documentation consists of one lexicon and some comparative work. No grammar was found, though a separate search might locate Remijsen's tone work. A curriculum would need new work with the Matbat community, including orthography development and handling of the tone system.
\end{document}